基于 CLIP 的零样本异常检测依赖固定的正常/异常文本原型，但固定原型缺少单张测试图的局部外观参照。仅用 CLIP 分数从测试图中选择 patch 会继承其局部误差；直接使用高频响应又会把规则纹理和结构边界误判为缺陷。本文提出\textbf{频谱可靠性引导的测试时语义原型校准}：在冻结的 CLIP patch 特征上计算边界感知的小波可靠性，用它温和调制正常/异常语义证据，构造逐图视觉原型，再以保守、非对称的方式校准固定文本原型。最终异常图仍由校准后的语义相似度产生，而非频率图融合。该框架把小波从异常评分器重新定位为逐图证据的可靠性估计器。为验证这一定位，我们预先规划纯原型路径与工程增强路径，比较直接图融合、频率特征增强、纯语义原型和完整方法，并通过预注册类别分组、正常图稳定性与非循环的漏检子集分析排除替代解释。当前稿件用于冻结实验判据，所有数字均为目标占位，正式版本将由真实评测替换。
